\documentclass{situationist-preprint}

\title{The Phantom Limb and the Body Schema}
\msid{PHE-14-b}
\msfamily{PHE-14}
\msdiscipline{Phenomenology}
\mstype{Companion manuscript}
\mskeywords{phantom limb; body schema; habitual body; zero-point; primacy of perception; Merleau-Ponty}

\begin{document}
\maketitle

\begin{abstract}
A person whose arm has been amputated may continue to feel it, to reach with it, and to locate pains in it, while knowing perfectly well that it is gone. Merleau-Ponty took the phantom limb as evidence that the body is given in a pre-reflective body schema that neither physiology nor psychology of representation explains. We argue that the case supports the primacy of perception defended in PHE-14-a. The phantom is a perceptual structure of the lived body as zero-point of orientation, which judgment does not override; and what does modify it, as mirror therapy shows, is a further perception. Knowledge of one's own body is founded on the schema, not the schema on knowledge.
\end{abstract}

\section{The case}

After amputation most patients report a phantom: the limb is felt in a definite posture, it may itch or hurt, and it can seem to move, sometimes voluntarily \citep{Ramachandran1998}. The patient does not believe that the limb is there. Asked, she says it is gone; yet she may reach for a cup with it, or step out of bed on a missing leg. Merleau-Ponty argued that neither of the explanations available to him fits. A physiological account, on which the phantom is the continued firing of nerve pathways, cannot explain why it responds to emotions and memories, and a psychological account, on which it is a representation or a refusal of mutilation, cannot explain why it persists against what the patient knows \citep[Part~I, ch.~1]{MerleauPonty1945}. He proposed instead that the phantom belongs to the \emph{habitual body}: the body as the system of possibilities for action that the patient has built up, which continues to orient her world, in contrast to the \emph{present body}, which has lost the limb.

\section{The argument}

\begin{definition}\label{d:schema}
The \emph{body schema} is the pre-reflective organization of the lived body as a field of possible movements, by which the body's parts are given as available for action and the perceptual field is oriented \citep{Gallagher2005}.
\end{definition}

The schema is distinct from the body image, the set of perceptions, beliefs and attitudes whose object is one's body \citep{Gallagher2005}. The phantom is a persistence of the schema, together with the experiences in which it shows itself, against the body image.

\begin{proposition}\label{p:zero}
The phantom limb belongs to the body as zero-point of orientation, not to the body as object.
\end{proposition}

\begin{proof}
The phantom is not given as an object in the sense of Definition~2.1 of PHE-08-a: it cannot be looked at, approached or explored from another side. It is given as the place from which a reach would start and as a part of the here relative to which things are near or far, which is how the body as zero-point is given (Definition~2.2 and Lemma~3.1 of PHE-08-a). The phantom therefore shows, in a case where the objective limb is absent, the organization that Theorem~3.2 of PHE-08-a describes: the orienting body is not exhausted by what of it can be given as an object.
\end{proof}

\begin{proposition}\label{p:founding}
The patient's knowledge that the limb is gone is founded on perception and does not modify the schema; what modifies the schema is a further perception.
\end{proposition}

\begin{proof}
The knowledge is gained by seeing and touching the stump, which are perceptual acts in the sense of Definition~2.2 of PHE-14-a; its evidence is founded on them in the sense of Definition~2.1 of PHE-14-a. That the knowledge does not modify the schema is the clinical datum. That a further perception does is shown by mirror therapy: when a mirror lets the patient see her intact hand in the place of the missing one, movements of the seen hand are felt in the phantom, and a phantom felt as clenched or paralysed can be relieved \citep{Ramachandran1996}. A visual perception changes what a judgment could not.
\end{proof}

The case thus displays both directions of the founding relation that PHE-14-a describes. The judgment about one's own body is founded on perception (Corollary~3.3 of PHE-14-a), and the perceptual layer of the lived body has an organization of its own, answerable to perception and not to the judgments founded on it.

\begin{objection}
Neuroscience explains the phantom by the persistence of a central representation of the body, a neuromatrix generating the experience of a body with or without input from the limbs \citep{Melzack1990}. The phenomenological description is superfluous.
\end{objection}

\begin{reply}
The neural account explains what causally sustains the schema; it does not compete with a claim about the founding of evidence. PHE-14-a concerns what the evidence of an act rests on for its subject, and the neuromatrix is no part of the patient's evidence about her body. The two accounts are about different relations, and the neural one is, moreover, itself founded on perception, in the sense of PHE-14-a, through the observations and records from which it was established.
\end{reply}

\section{Conclusion}

The phantom limb is a case in which a perceptual structure of the lived body persists against knowledge and yields to perception. It confirms the primacy of perception for the knowledge one has of one's own body and illustrates, in its starkest form, the claim of PHE-08-a that the body as zero-point is not given as an object.

\begin{thebibliography}{9}\small
\bibitem[Gallagher(2005)]{Gallagher2005} Gallagher, S. (2005). \emph{How the Body Shapes the Mind}. Oxford University Press.
\bibitem[Melzack(1990)]{Melzack1990} Melzack, R. (1990). Phantom limbs and the concept of a neuromatrix. \emph{Trends in Neurosciences}, 13(3), 88--92.
\bibitem[Merleau-Ponty(1945)]{MerleauPonty1945} Merleau-Ponty, M. (1945). \emph{Ph\'enom\'enologie de la perception}. Gallimard. Trans. D.~A. Landes, \emph{Phenomenology of Perception}, Routledge, 2012.
\bibitem[Ramachandran and Hirstein(1998)]{Ramachandran1998} Ramachandran, V.~S., and Hirstein, W. (1998). The perception of phantom limbs: The D.~O. Hebb lecture. \emph{Brain}, 121(9), 1603--1630.
\bibitem[Ramachandran and Rogers-Ramachandran(1996)]{Ramachandran1996} Ramachandran, V.~S., and Rogers-Ramachandran, D. (1996). Synaesthesia in phantom limbs induced with mirrors. \emph{Proceedings of the Royal Society of London B}, 263, 377--386.
\end{thebibliography}

\end{document}
